\documentclass[11pt]{article}
\usepackage[margin=1in]{geometry}
\usepackage{amsmath,amssymb}
\usepackage{booktabs}
\usepackage{enumitem}
\usepackage{xcolor}
\usepackage[colorlinks=true,linkcolor=blue!50!black,urlcolor=blue!50!black]{hyperref}
\usepackage{parskip}

\title{\vspace{-1.5em}SciNoNet: A Helmholtz-Potential Physics-Informed Network\\[-0.2em]
\large Architecture, Governing Physics, and Loss Functions}
\author{}
\date{}

% standardized (hatted) variables
\newcommand{\xh}{\hat{x}} \newcommand{\yh}{\hat{y}} \newcommand{\zh}{\hat{z}}
\newcommand{\tht}{\hat{t}}
\newcommand{\uh}{\hat{u}} \newcommand{\vh}{\hat{v}} \newcommand{\wh}{\hat{w}}
\newcommand{\phih}{\hat{\phi}} \newcommand{\psih}{\hat{\psi}}
\newcommand{\cph}{\hat{c}_p} \newcommand{\csh}{\hat{c}_s}
\newcommand{\dd}[2]{\frac{\partial #1}{\partial #2}}
\newcommand{\ddd}[2]{\frac{\partial^2 #1}{\partial #2^2}}

\begin{document}
\maketitle
\vspace{-3em}

\section{Overview}
The model reconstructs the three displacement components $u,v,w$ of guided
elastic (Lamb) waves on an aluminum plate. A neural network predicts four
\emph{Helmholtz potentials} from the space--time coordinates; the displacement is
obtained from the potentials by the Helmholtz decomposition, and the potentials
are constrained to satisfy two scalar wave equations (one longitudinal, one
shear), a gauge condition, and a rest initial condition. Every quantity is
expressed in standardized (non-dimensional) variables; Section~\ref{sec:std}
defines them and Table~\ref{tab:sym} lists every symbol with its value.

\section{Coordinates, fields, and parameters}
\label{sec:std}
The physical inputs are the position $(x,y,z)$ in mm and time $t$ in s; the
outputs are the displacements $(u,v,w)$ in mm. The plate is $300\times200\times2$
mm with the through-thickness direction $z$. The data are standardized to
zero-centered, $O(1)$ variables. With sample means $\mu_\bullet$ and standard
deviations as defined below,
\begin{equation}
\xh=\frac{x-\mu_x}{L},\quad \yh=\frac{y-\mu_y}{L},\quad
\zh=\frac{z-\mu_z}{L},\quad \tht=\frac{t-\mu_t}{s_t},
\label{eq:coordstd}
\end{equation}
where $L$ is a single in-plane length scale used for \emph{all three} spatial
axes (including $z$), and $s_t$ is the time scale. Using one length $L$ for $z$
as well (rather than the small plate thickness) keeps the spatial derivatives
balanced; this is captured by the ratio
\begin{equation}
\rho \equiv \frac{L}{L_z},\qquad L_z=L \;\Rightarrow\; \rho=1,
\end{equation}
which multiplies every $z$-derivative term below; the equations are written for
general $\rho$ and reduce to the implemented case at $\rho=1$. The fields are
standardized by a single common scale $s_f$ (their combined standard deviation):
\begin{equation}
\uh=\frac{u}{s_f},\quad \vh=\frac{v}{s_f},\quad \wh=\frac{w}{s_f}.
\label{eq:fieldstd}
\end{equation}

\begin{table}[h]
\centering\small
\begin{tabular}{lll}
\toprule
Symbol & Meaning & Value (this run) \\
\midrule
$\mu_x,\mu_y,\mu_z$ & coordinate means & data-dependent \\
$L$ & in-plane length scale ($=\!\mathrm{std}$ of $x,y$) & $\approx 89$ mm \\
$L_z$ & through-thickness scale ($=L$) & $\approx 89$ mm \\
$\rho=L/L_z$ & in-plane/thickness ratio & $1$ \\
$s_t$ & time scale ($=\!\mathrm{std}$ of $t$) & $\approx 1.73\times10^{-5}$ s \\
$s_f$ & field scale ($=\!\mathrm{std}$ of $u,v,w$) & $\approx 6.7\times10^{-6}$ mm \\
$c_p$ & longitudinal (P) wave speed & $6.3\times10^{6}$ mm/s \\
$c_s$ & shear (S) wave speed & $3.15\times10^{6}$ mm/s \\
$\cph=c_p s_t/L$ & standardized P speed & $\approx 1.22$ \\
$\csh=c_s s_t/L$ & standardized S speed & $\approx 0.61$ \\
$L_f$ & number of Fourier frequencies & $160$ \\
$\alpha$ & gradient-balancing target ratio & $0.3$ \\
$w_g,w_i$ & gauge / IC sub-weights & $1,\,1$ \\
$N_c,N_i$ & collocation / IC batch sizes & $2048,\,1024$ \\
$\Delta t$ & simulation timestep & $10^{-8}$ s \\
\bottomrule
\end{tabular}
\caption{Symbols and their values.}
\label{tab:sym}
\end{table}

\section{Network architecture}
\label{sec:arch}
The network maps the standardized coordinates $\mathbf{x}=(\xh,\yh,\zh,\tht)$ to
the four scalar potentials
\begin{equation}
\mathbf{p}(\mathbf{x}) = \big(\phih,\ \psih_x,\ \psih_y,\ \psih_z\big)
= \mathrm{MLP}\big(\gamma(\mathbf{x})\big),
\end{equation}
where $\phih$ is the (dimensionless) scalar potential and
$\boldsymbol{\psih}=(\psih_x,\psih_y,\psih_z)$ the vector potential.

\paragraph{Fourier-feature embedding.} The coordinates are first lifted by a
specialized random-Fourier embedding. With a frequency matrix
$\mathbf{B}\in\mathbb{R}^{4\times L_f}$ (columns $\mathbf{B}_{:,j}$) and a
per-frequency amplitude $A_j$,
\begin{equation}
\gamma(\mathbf{x}) = \Big[\,A_1\sin(2\pi\,\mathbf{x}\!\cdot\!\mathbf{B}_{:,1}),\dots,
A_{L_f}\sin(2\pi\,\mathbf{x}\!\cdot\!\mathbf{B}_{:,L_f}),\;
A_1\cos(\cdot),\dots,A_{L_f}\cos(\cdot)\,,\ \mathbf{x}\Big]\in\mathbb{R}^{2L_f+4}.
\label{eq:rff}
\end{equation}
The raw coordinates $\mathbf{x}$ are appended (the trailing term) so the network
can also represent slow, linear trends. Both $A_j$ and $\mathbf{B}$ are learnable;
$\mathbf{B}$ is initialized uniformly within and clamped after each step to
per-dimension bands taken from the wave dispersion curve. In standardized
cycles-per-unit (a physical frequency $f$ maps to $f\cdot\mathrm{std}$),
\begin{equation}
B_{x},B_{y},B_{z}\in[-k_{\max}L,\,k_{\max}L],\qquad B_{t}\in[0,\,f_{\max}s_t],
\end{equation}
with $k_{\max}=0.26$ cyc/mm ($=2\times0.13$, the band $\kappa\!\le\!0.8$ rad/mm)
and $f_{\max}=300$ kHz, giving $|B_{x,y,z}|\!\lesssim\!23$ and $B_t\!\lesssim\!5.2$.

\paragraph{MLP.} The embedding feeds a multilayer perceptron with three hidden
layers of width $256$ and $\tanh$ activations, ending in a linear layer with four
outputs (the potentials). The two wave speeds enter as fixed parameters
$\cph,\csh$ (Section~\ref{sec:phys}); they are stored as learnable
log-parameters but held fixed here so that $c_p/c_s$ stays at its physical value.

\section{Displacement from the potentials}
\label{sec:disp}
The Helmholtz decomposition writes the displacement field as the gradient of the
scalar potential plus the curl of the vector potential,
$\mathbf{u}=\nabla\Phi+\nabla\times\boldsymbol{\Psi}$, i.e.
\begin{equation}
u=\dd{\Phi}{x}+\dd{\Psi_z}{y}-\dd{\Psi_y}{z},\quad
v=\dd{\Phi}{y}+\dd{\Psi_x}{z}-\dd{\Psi_z}{x},\quad
w=\dd{\Phi}{z}+\dd{\Psi_y}{x}-\dd{\Psi_x}{y},
\label{eq:helm}
\end{equation}
where $\Phi,\boldsymbol{\Psi}$ are the \emph{physical} potentials. They are
related to the network's dimensionless outputs by a single scale $C=s_f L$,
\begin{equation}
\Phi=C\,\phih,\qquad \Psi_x=C\,\psih_x,\quad \Psi_y=C\,\psih_y,\quad \Psi_z=C\,\psih_z .
\end{equation}
From \eqref{eq:coordstd}, a physical derivative becomes a standardized one through
the chain rule, $\partial/\partial x=(1/L)\,\partial/\partial\xh$,
$\partial/\partial z=(1/L_z)\,\partial/\partial\zh$. Substituting into the first
line of \eqref{eq:helm},
\begin{align}
u &= C\Big(\tfrac{1}{L}\phih_{\xh}+\tfrac{1}{L}\psih_{z,\yh}-\tfrac{1}{L_z}\psih_{y,\zh}\Big)
= \frac{C}{L}\Big(\phih_{\xh}+\psih_{z,\yh}-\tfrac{L}{L_z}\psih_{y,\zh}\Big) \notag\\
&= s_f\big(\phih_{\xh}+\psih_{z,\yh}-\rho\,\psih_{y,\zh}\big),
\end{align}
using $C/L=s_f$. Dividing by $s_f$ gives the standardized displacement the data
loss compares against (subscripts after a comma denote partial derivatives):
\begin{equation}
\boxed{\;
\uh=\phih_{\xh}+\psih_{z,\yh}-\rho\,\psih_{y,\zh},\quad
\vh=\phih_{\yh}+\rho\,\psih_{x,\zh}-\psih_{z,\xh},\quad
\wh=\rho\,\phih_{\zh}+\psih_{y,\xh}-\psih_{x,\yh}.\;}
\label{eq:disp_std}
\end{equation}
All derivatives are obtained by automatic differentiation of the network output.

\section{Governing physics}
\label{sec:phys}

\subsection{Wave equations and their standardized residual}
For an isotropic elastic solid the scalar potential carries the longitudinal (P)
wave at speed $c_p$ and each component of the vector potential carries the shear
(S) wave at speed $c_s$:
\begin{equation}
\ddd{\Phi}{t}=c_p^2\,\nabla^2\Phi,\qquad
\ddd{\Psi_k}{t}=c_s^2\,\nabla^2\Psi_k\quad(k=x,y,z),
\label{eq:wave_phys}
\end{equation}
with $\nabla^2=\partial_{xx}+\partial_{yy}+\partial_{zz}$. The two fundamental
Lamb modes (symmetric $S_0$, antisymmetric $A_0$) arise from the superposition of
$\Phi$ and $\boldsymbol{\Psi}$ through the plate thickness. We derive the
residual for a generic potential $q\in\{\phih,\psih_x,\psih_y,\psih_z\}$ with
speed $c\in\{c_p,c_s\}$. Writing the physical potential $Q=C\,q$ and using the
chain rule on \eqref{eq:wave_phys},
\begin{align}
\ddd{Q}{t}=\frac{C}{s_t^2}\,q_{\tht\tht},\qquad
\nabla^2 Q = C\Big[\tfrac{1}{L^2}(q_{\xh\xh}+q_{\yh\yh})+\tfrac{1}{L_z^2}q_{\zh\zh}\Big].
\end{align}
Substituting and cancelling $C$,
\begin{equation}
\frac{1}{s_t^2}q_{\tht\tht}
= c^2\Big[\tfrac{1}{L^2}(q_{\xh\xh}+q_{\yh\yh})+\tfrac{1}{L_z^2}q_{\zh\zh}\Big]
\;\Longrightarrow\;
q_{\tht\tht}=\hat c^2\,(q_{\xh\xh}+q_{\yh\yh})+\hat c^2\rho^2\,q_{\zh\zh},
\end{equation}
where $\hat c=c\,s_t/L$ is the standardized speed and we used
$c^2 s_t^2/L^2=\hat c^2$ and $c^2 s_t^2/L_z^2=\hat c^2\rho^2$ (since
$1/L_z^2=\rho^2/L^2$). Dividing through by $\hat c^2\rho^2$ to keep all terms
$O(1)$ gives the residual used in the loss,
\begin{equation}
\boxed{\;
r_q \;=\; \frac{1}{\hat c^2\rho^2}\,q_{\tht\tht}
\;-\;\frac{1}{\rho^2}\big(q_{\xh\xh}+q_{\yh\yh}\big)\;-\;q_{\zh\zh}\;,}
\label{eq:res}
\end{equation}
with $\hat c=\cph$ for $q=\phih$ and $\hat c=\csh$ for the three
$\psih_k$. At $\rho=1$ this is $r_q=q_{\tht\tht}/\hat c^2-(q_{\xh\xh}+q_{\yh\yh}+q_{\zh\zh})$.
A solution of the wave equation makes $r_q=0$.

\subsection{Coulomb gauge}
The vector potential is not unique: $\boldsymbol{\Psi}\!\to\!\boldsymbol{\Psi}+\nabla\chi$
changes nothing in $\nabla\times\boldsymbol{\Psi}$, hence nothing in $\mathbf u$.
This freedom is removed by the Coulomb gauge $\nabla\!\cdot\!\boldsymbol{\Psi}=0$.
In standardized variables,
\begin{equation}
\nabla\!\cdot\!\boldsymbol{\Psi}
=\dd{\Psi_x}{x}+\dd{\Psi_y}{y}+\dd{\Psi_z}{z}
=\frac{C}{L}\Big(\psih_{x,\xh}+\psih_{y,\yh}+\rho\,\psih_{z,\zh}\Big),
\end{equation}
so the (scale-free) gauge residual is
\begin{equation}
\boxed{\;g=\psih_{x,\xh}+\psih_{y,\yh}+\rho\,\psih_{z,\zh}\;.}
\label{eq:gauge}
\end{equation}

\subsection{Initial condition}
The plate is at rest before excitation: at $t=0$ the displacement and velocity
vanish, which (with the gauge) is enforced by requiring the potentials and their
time-derivatives to vanish on the $t=0$ plane,
\begin{equation}
\mathbf{p}\big|_{\tht_0}=\mathbf 0,\qquad
\dd{\mathbf{p}}{\tht}\Big|_{\tht_0}=\mathbf 0,\qquad
\tht_0=\frac{0-\mu_t}{s_t}.
\label{eq:ic}
\end{equation}
The eight scalar residuals are stacked into
$\mathbf{r}_{\mathrm{ic}}=(\,\mathbf{p},\ \partial\mathbf{p}/\partial\tht\,)\in\mathbb{R}^{8}$.

\section{Loss functions}
\label{sec:loss}
Let $\{\mathbf{x}_n\}_{n=1}^{N}$ be the data points (with measured displacements
$\hat{\mathbf u}^{\mathrm{data}}_n$), $\{\mathbf{x}^c_i\}_{i=1}^{N_c}$ collocation
points sampled uniformly over the standardized domain box for the wave and gauge
residuals, and $\{\mathbf{x}^{\mathrm{ic}}_i\}_{i=1}^{N_i}$ points on the $t=0$
plane. Each loss is a mean square.

\paragraph{Data loss.} Compares the displacement \eqref{eq:disp_std} (derived from
the predicted potentials) to the measurements:
\begin{equation}
\mathcal{L}_{\mathrm{data}}=\frac{1}{N}\sum_{n=1}^{N}
\Big[(\uh-\uh^{\mathrm{data}})_n^2+(\vh-\vh^{\mathrm{data}})_n^2+(\wh-\wh^{\mathrm{data}})_n^2\Big].
\end{equation}

\paragraph{Wave-residual loss.} Sums the squared residual \eqref{eq:res} over the
four potentials:
\begin{equation}
\mathcal{L}_{\mathrm{wave}}=\frac{1}{N_c}\sum_{i=1}^{N_c}\sum_{q\in\{\phih,\psih_x,\psih_y,\psih_z\}} r_q(\mathbf{x}^c_i)^2 .
\end{equation}

\paragraph{Gauge loss.} Penalizes departures from $\nabla\!\cdot\!\boldsymbol{\Psi}=0$ via \eqref{eq:gauge}:
\begin{equation}
\mathcal{L}_{\mathrm{gauge}}=\frac{1}{N_c}\sum_{i=1}^{N_c} g(\mathbf{x}^c_i)^2 .
\end{equation}

\paragraph{Initial-condition loss.} Enforces \eqref{eq:ic}:
\begin{equation}
\mathcal{L}_{\mathrm{ic}}=\frac{1}{N_i}\sum_{i=1}^{N_i}\big\lVert \mathbf{r}_{\mathrm{ic}}(\mathbf{x}^{\mathrm{ic}}_i)\big\rVert^2
=\frac{1}{N_i}\sum_{i=1}^{N_i}\sum_{q}\Big(q^2+q_{\tht}^2\Big)\Big|_{\tht_0}.
\end{equation}

\paragraph{Total loss and gradient-norm balancing.} The physics terms are grouped as
\begin{equation}
\mathcal{L}_{\mathrm{phys}}=\mathcal{L}_{\mathrm{wave}}+w_g\,\mathcal{L}_{\mathrm{gauge}}+w_i\,\mathcal{L}_{\mathrm{ic}},
\qquad
\mathcal{L}=\mathcal{L}_{\mathrm{data}}+\lambda\,\mathcal{L}_{\mathrm{phys}}.
\end{equation}
The physics weight $\lambda$ is set each epoch so that the physics gradient is a
fixed fraction $\alpha$ of the data gradient (preventing the large second-order
physics gradients from overwhelming the data fit),
\begin{equation}
\lambda^{\star}=\alpha\,\frac{\big\lVert\nabla_\theta\mathcal{L}_{\mathrm{data}}\big\rVert}
{\big\lVert\nabla_\theta\mathcal{L}_{\mathrm{phys}}\big\rVert},
\qquad
\lambda\leftarrow\beta\,\lambda+(1-\beta)\,\lambda^{\star},
\end{equation}
where $\theta$ are the network weights, $\nabla_\theta$ the gradient with respect
to them, $\lVert\cdot\rVert$ the Euclidean norm, and $\beta=0.9$ an
exponential-moving-average factor. A short data-only warmup ($\lambda=0$ for the
first $12$ epochs) precedes the physics phase.

\section{Training}
Optimizer: Adam, learning rate $2\times10^{-3}$, gradient-norm clipping at $1$,
and a plateau learning-rate schedule. All derivatives (first order for the
displacement, second order for the wave residual) use automatic differentiation
in double precision. The best model is the one with the lowest validation data
loss. Data and collocation are drawn in minibatches (data batch $16384$).

\end{document}
